\documentclass[10pt]{article}
\usepackage[T1]{fontenc}
\usepackage[utf8]{inputenc}
\usepackage[french]{babel}
\usepackage{amsmath,amssymb,mathrsfs,mathpazo}
\usepackage{geometry}
\geometry{a4paper,margin=2cm}
\usepackage{graphicx,array,tabularx,tabularray,colortbl,multirow,multicol,booktabs,enumitem,pifont,ulem,fancybox,diagbox,cancel}
\usepackage{pstricks,pstricks-add,pst-plot,pst-tree,pst-node,pst-text,pst-func,pst-math,pst-bspline,pst-3dplot,pst-eucl,pst-solides3d}
\let\pscalculate\undefined
\usepackage{pst-bezier}
\usepackage{tikz,pgfplots,tkz-tab}
\usetikzlibrary{arrows,arrows.meta,patterns,calc,positioning,trees,decorations.markings,shapes.geometric,shapes.misc,shapes.arrows,intersections,decorations.pathmorphing,decorations.pathreplacing,backgrounds,shadows,babel}
\pgfplotsset{compat=1.18}
\usepgfplotslibrary{fillbetween}
\usepackage[np]{numprint}
\usepackage{listings,makecell,pas-tableur,twoopt}
\let\barre\undefined
\usepackage[pstricks]{tablvar}
\definecolor{darkgreen}{RGB}{0,100,0}
\providecommand{\red}{\color{red}}
\providecommand{\blue}{\color{blue}}
\providecommand{\green}{\color{green}}
\providecommand{\black}{\color{black}}
\providecommand{\white}{\color{white}}
\providecommand{\gray}{\color{gray}}
\providecommand{\cyan}{\color{cyan}}
\providecommand{\magenta}{\color{magenta}}
\providecommand{\yellow}{\color{yellow}}
\providecommand{\violet}{\color{violet}}
\providecommand{\orange}{\color{orange}}
\providecommand{\euro}{\texteuro}
\providecommand{\EUR}{\texteuro}
\providecommand{\ieme}{\textsuperscript{e}}
\providecommand{\degres}{\ensuremath{^\circ}}
\providecommand{\degree}{\ensuremath{^\circ}}
\providecommand{\vv}[1]{\overrightarrow{#1}}
\providecommand{\Coord}[1]{#1}
\providecommand{\decosix}{}
\providecommand{\starredbullet}{\ensuremath{\bullet}}
\pagestyle{empty}
\setlength{\parindent}{0pt}
\usepackage{adjustbox}\DeclareUnicodeCharacter{2236}{:}\DeclareUnicodeCharacter{2190}{\ensuremath{\leftarrow}}
\providecommand{\R}{\mathbb{R}}
\providecommand{\N}{\mathbb{N}}
\providecommand{\Z}{\mathbb{Z}}
\providecommand{\Q}{\mathbb{Q}}
\providecommand{\D}{\mathbb{D}}
\providecommand{\C}{\mathbb{C}}
\providecommand{\e}{{\mathrm{e}}}
\providecommand{\vect}[1]{\overrightarrow{#1}}
\providecommand{\vectt}[1]{\overrightarrow{\mathrm{#1}}}
\providecommand{\vvt}[1]{\overrightarrow{\mathrm{#1}}}
\providecommand{\np}[2][]{{#2~#1}}
\providecommand{\vv}[1]{{\overrightarrow{#1}}}
\renewcommand{\d}{{\mathrm{d}}}
\providecommand{\degres}{{^\circ}}
\providecommand{\degre}{{^\circ}}
\providecommand{\blue}{\color{blue}}
\providecommand{\red}{\color{red}}
\providecommand{\enskip}{\quad}
\providecommand{\llbracket}{\lbrack\!\lbrack}
\providecommand{\rrbracket}{\rbrack\!\rbrack}
\providecommand{\euro}{{\text{€}}}
\providecommand{\no}{{\text{n°}}}
\providecommand{\interval}[2]{{[#1;#2]}}
\providecommand{\intervFO}[2]{{[#1;#2[}}
\providecommand{\dotfill}{\ldots}

\providecommand{\ds}{\displaystyle}
\providecommand{\textding}[1]{#1}
\providecommand{\Oij}{\(\mathrm{O};\vec{\imath},\vec{\jmath}\)}
\providecommand{\Oijk}{\(\mathrm{O};\vec{\imath},\vec{\jmath},\vec{k}\)}
\providecommand{\Ouv}{\(\mathrm{O};\vec{u},\vec{v}\)}
\providecommand{\intervFF}[2]{\ensuremath{[#1;#2]}}
\renewcommand{\euro}{\ifmmode\text{\texteuro}\else\texteuro\fi}
\DeclareUnicodeCharacter{00B0}{\ifmmode^{\circ}\else\textdegree\fi}
\DeclareUnicodeCharacter{20AC}{\ifmmode\text{\texteuro}\else\texteuro\fi}
% Énoncés : Éducation nationale. Transcription et corrigés : APMEP, Denis Vergès et contributeurs. Découpage et classement : CoopMaths, Sébastien Lozano, Éric Elter et contributeurs. Adaptation de la présentation : Maths Explorer.
% CC BY-SA : https://coopmaths.fr/www/about/
% https://coopmaths.fr/alea/static/bac/2021/tex/bac_2021_03_sujet2_4.tex
\begin{document}
\noindent\textit{APMEP / CoopMaths. CC BY-SA. Présentation adaptée par Maths Explorer.}\par\medskip
\medskip

\begin{tabularx}{\linewidth}{|X|}\hline
\textbf{Principaux domaines abordés: Fonction exponentielle; dérivation.}\\ \hline
\end{tabularx}

\bigskip

\parbox{0.35\linewidth}{Le graphique ci-contre représente, dans un repère orthogonal, les courbes 
$\mathcal{C}_f$ et $\mathcal{C}_g$ des fonctions $f$ et $g$ définies sur $\R$ par :

\[f(x) = x^2\text{e}^{-x}\quad \text{et} \quad g(x) = \text{e}^{-x}.\]}
\hfill \parbox{0.64\linewidth}{\psset{xunit=1.8cm,yunit=0.9cm}
\begin{pspicture*}(-2.5,-1)(2.5,9)
\multido{\n=0+2}{5}{\psline[linewidth=0.1pt](-2.5,\n)(2.5,\n)}
\multido{\n=-2+1}{5}{\psline[linewidth=0.1pt](\n,-1)(\n,9)}
\psaxes[linewidth=1.25pt,labelFontSize=\scriptstyle,Dy=2]{->}(0,0)(-2.5,-1)(2.5,9)
\psplot[linewidth=1.25pt,linecolor=blue,plotpoints=2000]{-2.5}{2.5}{x dup mul 2.71828 x exp div}
\psplot[linewidth=1.25pt,linecolor=red,plotpoints=2000]{-2.5}{2.5}{2.71828 x neg exp}
\psline[linestyle=dashed,linewidth=1.25pt](-0.6,-1)(-0.6,9)
\psdots(-0.6,0.656)(-0.6,1.822)
\uput[r](-0.6,0.656){\footnotesize $N$}\uput[r](-0.6,1.822){\footnotesize $M$}
\uput[r](-1.4,7){\blue $\mathcal{C}_f$}\uput[l](-2.1,7.4){\red $\mathcal{C}_g$}
\end{pspicture*}}

\textbf{La question 3 est indépendante des questions 1 et 2.}

\medskip

\begin{enumerate}
\item 
	\begin{enumerate}
		\item Déterminer les coordonnées des points d'intersection de $\mathcal{C}_f$ et $\mathcal{C}_g$.
		\item Étudier la position relative des courbes $\mathcal{C}_f$ et $\mathcal{C}_g$.
	\end{enumerate}
\item  Pour tout nombre réel $x$ de l'intervalle $[-1~;~1]$, on considère les points $M$ de coordonnées $(x~;~f(x))$ et $N$ de coordonnées $(x~;~g(x))$, et on note $d(x)$ la distance $MN$. On admet que : $d(x)= \text{e}^{-x} - x^2\text{e}^{-x}$.

On admet que la fonction $d$ est dérivable sur l'intervalle $[-1~;~1]$ et on note $d'$ sa fonction dérivée.
	\begin{enumerate}
		\item Montrer que $d'(x) = \text{e}^{-x}\left(x^2 - 2x - 1\right)$. 
		\item En déduire les variations de la fonction $d$ sur l'intervalle $[-1~;~1]$.
		\item Déterminer l'abscisse commune $x_0$ des points $M_0$ et $N_0$ permettant d'obtenir une
distance $d\left(x_0\right)$ maximale, et donner une valeur approchée à $0,1$ près de la distance $M_0N_0$.
	\end{enumerate}
\item  Soit $\Delta$ la droite d'équation $y = x + 2$.

On considère la fonction $h$ dérivable sur $\R$ et définie par: $h(x) = \text{e}^{-x} - x - 2$.

En étudiant le nombre de solutions de l'équation $h(x) = 0$, déterminer le nombre de points d'intersection de la droite $\Delta$ et de la courbe $\mathcal{C}_g$.
\end{enumerate}

\end{document}
